\chapter{Methodology}

\section{Exploratory data analysis}
EDA was used to determine what could be measured fairly before selecting a model. For RIRD,
we measured restaurant review counts, review length, sentence count, rating distribution and
the number of relevant restaurants per query. The 50 restaurants have between 406 and 2,281
reviews, so an unconstrained top-$k$ aggregation could favour restaurants with more
opportunities to obtain a high score. RIRD also has 100 queries and five-annotator judgements,
making query-level paired evaluation possible.

Wongnai contains real Thai review text and star ratings but no restaurant identifier. This EDA
result changed the claim of the project: Wongnai can support a sentence-level Thai taste study,
not Thai restaurant ranking. Approximately 10,000 reviews were sampled for the study, yielding
93,810 sentences.

\section{Pre-processing}
English reviews are split at sentence-ending punctuation and newlines. Empty units are removed,
while review and restaurant IDs are retained so every retrieved sentence can be traced to its
source. Thai reviews are segmented with PyThaiNLP because Thai words are not separated
consistently by spaces. BM25 uses lower-case English
tokens and Thai tokenisation. Dense embeddings use
\texttt{paraphrase-multilingual-mpnet-base-v2}, following the sentence-embedding approach of
\citet{reimers2019sbert}.

For the live English path, the full 240,660-sentence corpus is reconstructed in the same order
as the cached embedding matrix. Sentences shorter than ten words are excluded from live
candidate generation because inspection found generic fragments such as ``Solid meal!'' to be
poor explanations. This filtering is applied only to the live evidence path; benchmark results
are computed from the controlled exported evaluation data.

\section{Baseline models and aggregation}
The primary lexical baseline is BM25 \citep{robertson2009bm25}. The dense baseline encodes
each review or sentence and computes cosine similarity with the query. We evaluate both at
review and sentence level. For a query, each restaurant receives the scores of all its units,
then we apply:
\begin{align}
 f_{\max}(S_r)&=\max S_r,\\
 f_{\mathrm{mean}}(S_r)&=\frac{1}{|S_r|}\sum_{s\in S_r}s,\\
 f_{\mathrm{top}k}(S_r)&=\frac{1}{k}\sum_{s\in \operatorname{TopK}(S_r)}s.
\end{align}
We test $k\in\{3,5,10\}$. This shared-unit design is the key methodological control: BM25 is
not given a concatenated restaurant document while dense retrieval is given pooled reviews.

\section{Proposed model and system architecture}
The proposed evaluated model is dense sentence retrieval with top-10 mean aggregation. The
live demo adds query rewriting and selective cross-encoder reranking:
\begin{enumerate}
  \item Groq rewrites vague phrases and converts exclusion language into reviewer-oriented
  attributes;
  \item the multilingual bi-encoder retrieves candidate sentences;
  \item a cross-encoder reranks candidates only for explicit negation or exclusion;
  \item top-ten scores are averaged and the five highest original sentences are shown.
\end{enumerate}

\begin{figure}[H]
\centering
\begin{tikzpicture}[
node distance=6mm and 7mm,
block/.style={draw,rounded corners,align=center,minimum height=10mm,text width=2.8cm,fill=blue!7},
decision/.style={draw,diamond,align=center,aspect=2,fill=orange!12},
arrow/.style={-{Latex[length=2mm]},thick}]
\node[block] (q) {User query\\English or Thai};
\node[block,right=of q] (rw) {LLM rewrite\\vague/negative language};
\node[block,right=of rw] (be) {BM25 or\\multilingual bi-encoder};
\node[decision,below=of rw] (d) {Explicit\\exclusion?};
\node[block,right=of be] (ag) {Top-$k$ mean\\restaurant score};
\node[block,below=of ag] (ce) {Selective cross-encoder\\top 20 restaurants};
\node[block,left=of ce] (ev) {Original review\\sentences};
\draw[arrow](q)--(rw); \draw[arrow](rw)--(be); \draw[arrow](be)--(ag);
\draw[arrow](rw)--(d); \draw[arrow](d)--node[right,font=\scriptsize]{yes}(ce);
\draw[arrow](ce)--(ag); \draw[arrow](ag)--(ev);
\end{tikzpicture}
\caption{Architecture of the evaluated retrieval system and the improved live demo.}
\label{fig:architecture}
\end{figure}

Plain positive queries use the dense baseline. Queries containing patterns such as ``no'',
``not'', ``without'', ``avoid'' or ``instead of'' use the reranker on the top 20 restaurants and
15 candidate sentences per restaurant. This reduces a naive 1,500 candidate comparisons to
approximately 300 while preserving the stronger baseline for ordinary queries.

\section{Explanation extraction}
For each restaurant, the top ten unit scores are averaged. The five highest-scoring original
sentences are returned as explanations. The system never generates synthetic evidence. The
evidence CSV stores query, restaurant ID, rank, score, source review ID, sentence text and
random seed, allowing the output to be reproduced and inspected.

\section{Evaluation design}
RIRD ranking is measured with nDCG@5, nDCG@10, MAP, MRR, Recall@5/10 and Precision@5/10.
Confidence intervals use query-level bootstrap resampling and model differences use paired
randomisation tests with Holm correction. Review-volume bias is measured by Pearson and
Spearman correlations between restaurant score and review count before and after capping every
restaurant at 400 reviews. Explanations are assessed by shown-versus-random deletion and human
support labels on a 0--2 scale. Thai retrieval uses a blind pool from BM25 and dense top
results, with query-level nDCG@10 and precision@10 plus annotator agreement.
